%------------------------------------------------------------------------------%
\begin{question}
  Considerando a equação
  \[
    x^2-4x+4=0
  \]
  \vspace{-\baselineskip}
  \begin{enumerate}
    \item identifique a curva
    \item determine seus principais valores e pontos
    \item esboce a curva
  \end{enumerate}
\end{question}

%------------------------------------------------------------------------------%
 \begin{resolution}{Classificação}
   Completando o quadrado

   \begin{align*}
     x^2-4x+4&=0
     \\
     (x-2)^2&=0
     \\
     x&=2
   \end{align*}

   A equação representa uma única reta

   \[
     x=2
   \]

   Entretanto, como a equação original pode ser escrita como

   \[
     (x-2)(x-2)=0
   \]

   essa reta aparece com multiplicidade dois

   Portanto, trata-se de uma cônica degenerada formada
   por uma reta dupla paralela ao eixo $y$
 \end{resolution}

%------------------------------------------------------------------------------%
\begin{resolution}{Esboço}
  \centering
  \includegraphics{I-07-Identificacao-ex-G}
\end{resolution}

%------------------------------------------------------------------------------%
